\section{One cube and the continuum gap}\label{sec:cube}

\emph{Setting.} Let $G\subset\U(N)$ be compact and nontrivial and $\varphi$ a
positive continuous class function on $G$, such as $e^{N\beta\ReTr(\cdot)}$. On
the torus $\Lambda_L=(\Z/L\Z)^4$, $L$ even and $L\ge2n+4$, with links $E$ and
plaquettes $P$, let
\[
\mu(dU)=\frac1Z\prod_{p\in P}\varphi(U_p)\prod_{e\in E}dU_e .
\]
A link $(y,\nu)$ has base point $y$; differences of base points are taken in
$(-L/2,L/2]^4$ and $|\cdot|_\infty$ is the corresponding distance. For
$x\in\Lambda_L$ let $V_x$ be the set of links with base point in $x+[0,n)^4$;
every link lies in exactly $n^4$ of the $L^4$ cubes $V_x$. For $V\subset E$ let
$\partial V$ be the links outside $V$ that share a plaquette with a link of
$V$; the conditional law $\gamma^0_V(\cdot\mid U)$ of $U_V$ given the other links
has density proportional to $\prod_{p\cap V\ne\emptyset}\varphi(U_p)$ and depends
only on $U_{\partial V}$. Let $\mathcal P_xf(U)=\int f\,d\gamma^0_{V_x}(\cdot\mid U)$,
the conditional expectation given the links outside $V_x$, an orthogonal
projection on $L^2(\mu)$, and let $\mathcal A$ be the closed subspace of
gauge-invariant functions. By gauge covariance of $\gamma^0_{V_x}$, $\mathcal P_x$
maps $\mathcal A$ into itself.

\begin{definition}\label{def:bp}
For $n\ge2$ and $\varpi>0$, the condition $\BP(n,\varpi)$ for $\mu$ is
\[
\varpi\Var_\mu(f)\le\frac1{n^4}\sum_{x\in\Lambda_L}\mu\bigl((f-\mathcal P_xf)^2\bigr)
\qquad\text{for every real }f\in\mathcal A .
\]
\end{definition}

Put $q_n=(2+1/n)^4$, so that $16<q_n\le\frac{625}{16}$ for $n\ge2$, and let
$\mathcal K=L^{-4}\sum_x\mathcal P_x$ be the chain that resamples one uniformly
chosen cube; for $Y\subset\Lambda_L$ let $\mathcal K_Y=L^{-4}(\sum_{x\in Y}\mathcal P_x
+(L^4-|Y|)I)$ be the same chain with the moves outside $Y$ suppressed. Write
$y\rightsquigarrow z$ if $\partial V_z\cap V_y\ne\emptyset$. A link of $\partial
V_z$ has its base point in $z+[-1,n]^4$, so $y\rightsquigarrow z$ implies
$|z-y|_\infty\le n$, and for each $z$ at most $(2n+1)^4=q_nn^4$ points $y$ have
$y\rightsquigarrow z$. For finite $S\subset E$ let $Y_1(S)=\{y:V_y\cap S\ne
\emptyset\}$ and $Y_{j+1}(S)=Y_j(S)\cup\{y:y\rightsquigarrow z\text{ for some }
z\in Y_j(S)\}$, $Y_0(S)=\emptyset$.

\begin{lemma}[finite speed]\label{lem:speed}
Let $f$ be a bounded function of the links of a finite set $S$, let $m_0\ge0$
and $Y=Y_{m_0}(S)$. Then $\mathcal K_Y^tf$ depends only on the links of
$S\cup\bigcup_{y\in Y}\partial V_y$, and for every integer $t\ge0$
\[
\bigl\|\mathcal K^tf-\mathcal K_Y^tf\bigr\|_\infty\le\frac{2|S|}{q_n}\,\|f\|_\infty
\sum_{m>m_0}\frac{(q_n\tau)^m}{m!},\qquad\tau=\frac{tn^4}{L^4}.
\]
\end{lemma}

\begin{proof}
If $h$ depends only on the links of $W$, then $\mathcal P_xh=h$ when $V_x\cap W=
\emptyset$, and otherwise $\mathcal P_xh$ depends only on the links of $(W\setminus
V_x)\cup\partial V_x$; induction gives the first claim. For the second run two
chains from the same configuration with the same uniformly chosen cubes
$x_1,x_2,\dots$. If $x_s\notin Y$ the first chain resamples $V_{x_s}$ and the
second does nothing; if $x_s\in Y$ and the two configurations agree on
$\partial V_{x_s}$, both receive one common sample, and otherwise they resample
independently. The marginals are the chains $\mathcal K$ and $\mathcal K_Y$. A
link where the configurations differ at time $t$ lies in $V_{x_{s_m}}$ for some
times $s_1<\dots<s_m\le t$ with $x_{s_1}\notin Y$ and $x_{s_l}\rightsquigarrow
x_{s_{l+1}}$, by induction on $t$. If the set of such links meets $S$, then
$x_{s_m}\in Y_1(S)$ and $x_{s_{m-l}}\in Y_{l+1}(S)$, so $m>m_0$. For fixed $m$
there are at most $t^m/m!$ choices of times, $n^4|S|$ of $x_{s_m}$ and $q_nn^4$
of each earlier cube, each sequence having probability $L^{-4m}$; the
difference of $f$ is at most $2\|f\|_\infty$.
\end{proof}

For finite $S,S'\subset E$ let $d(S,S')$ be the least $|b-b'|_\infty$ over base
points $b$ of links of $S$ and $b'$ of links of $S'$.

\begin{theorem}[decay from one cube]\label{thm:decay}
Suppose $\BP(n,\varpi)$ holds for $\mu$. Then $\varpi\le4$, and for bounded
real $f,g\in\mathcal A$ depending on the links of finite sets $S_f,S_g$ with
$d=d(S_f,S_g)\ge3n-2$,
\[
|\Cov_\mu(f,g)|\le\Bigl(e+\frac{k|S_f||S_g|}{64e^2}\Bigr)\|f\|_\infty\|g\|_\infty\,
e^{-\varpi k/(e^2q_n)},\qquad k=\Bigl\lfloor\frac{d+2}n\Bigr\rfloor-2 .
\]
In particular $|\Cov_\mu(f,g)|\le C(1+d)e^{-\varpi d/(e^2q_nn)}$ for every $d$,
with $C$ depending only on $|S_f|,|S_g|,\|f\|_\infty,\|g\|_\infty$.
\end{theorem}

\begin{proof}
\emph{The bound $\varpi\le4$.} Take $f=\ReTr U_p$ for one plaquette. Each
$I-\mathcal P_x$ is an orthogonal projection annihilating constants, so
$\mu((f-\mathcal P_xf)^2)\le\Var_\mu(f)$, and it vanishes unless $V_x$ contains
a link of $p$, which happens for at most $4n^4$ points $x$.

\emph{Spectral part.} On the functions of $\mathcal A$ with mean zero, $\mathcal K$
is self-adjoint with $0\le\mathcal K\le I$, and $\BP(n,\varpi)$ reads
$\ip h{(I-\mathcal K)h}\ge\epsilon\|h\|^2$ with $\epsilon=\varpi n^4/L^4$. So
$\|\mathcal K^Th\|_2\le e^{-\epsilon T}\|h\|_2$, and from
$f-\mu(f)=\mathcal K^T(f-\mu(f))+\sum_{t<T}\mathcal K^t(I-\mathcal K)f$,
\[
\Cov_\mu(f,g)=\ip{\mathcal K^T(f-\mu f)}{g-\mu g}+\sum_{t=0}^{T-1}
\ip{\mathcal K^tf}{(I-\mathcal K)g},
\]
whose first term is at most $e^{-\epsilon T}\|f\|_\infty\|g\|_\infty$.

\emph{Local part.} Let $Z_g=Y_1(S_g)$, so $|Z_g|\le n^4|S_g|$ and $(I-\mathcal K)g=
L^{-4}\sum_{x\in Z_g}(g-\mathcal P_xg)$; if $h$ does not depend on the links of
$V_x$ then $\ip h{g-\mathcal P_xg}=0$. Put $m_0=k-1$ and $Y=Y_{m_0}(S_f)$. A link
of $V_x$, $x\in Z_g$, lying in $S_f$ or in $\partial V_y$ with $y\in Y$ would join
base points of $S_f$ and $S_g$ at distance at most $n(m_0+3)-3\le n\lfloor
(d+2)/n\rfloor-3<d$, which is impossible. So $\mathcal K_Y^tf$ does not depend on
the links of $V_x$, $x\in Z_g$, and by Lemma~\ref{lem:speed}
\[
\bigl|\ip{\mathcal K^tf}{(I-\mathcal K)g}\bigr|=\bigl|\ip{\mathcal K^tf-\mathcal K_Y^tf}
{(I-\mathcal K)g}\bigr|\le\frac{4|S_f||S_g|n^4}{q_nL^4}\|f\|_\infty\|g\|_\infty
\sum_{m\ge k}\frac{(q_n\tau_t)^m}{m!}.
\]
Summing over $t<T$ with $\tau=Tn^4/L^4$, the second term is at most
\[
\frac{4|S_f||S_g|\tau}{q_n}\,\|f\|_\infty\|g\|_\infty\sum_{m\ge k}
\frac{(q_n\tau)^m}{m!}.
\]

\emph{Choice of $T$.} Take $T=\lfloor kL^4/(e^2q_nn^4)\rfloor$, so that $z=q_n
\tau\le k/e^2$ and $\epsilon T\ge\varpi k/(e^2q_n)-1$, using $\varpi\le4$ and
$L^4\ge4n^4$. Since $k!\ge(k/e)^k$, $\sum_{m\ge k}z^m/m!\le(ez/k)^ke^z\le
e^{-k(1-e^{-2})}$, and $\tau/q_n\le k/(256e^2)$. As $\varpi/(e^2q_n)<1-e^{-2}$,
both terms are bounded as stated. For the second form use $k>d/n-3$, and for
$d<3n-2$ the trivial bound $\|f\|_\infty\|g\|_\infty$.
\end{proof}

A \emph{torus limit} at $\beta$ is a state $\omega$ on the gauge-invariant local
functions of $\Z^4$ with $\omega(O)=\lim_j\mu_{L_j}(O)$ along some $L_j\to\infty$;
by compactness every sequence has a subsequence along which one exists.

\begin{corollary}[lattice gap]\label{cor:latticegap}
Let $\beta\ge0$ and suppose $\BP(n,\varpi)$ holds for $\mu_L$ for all $L$ in an
infinite set of even integers $L\ge2n+4$. Then every torus limit along this
set is invariant under translations, reflection positive, and its transfer
matrix has spectral gap at least
\[
m_n=\frac{\varpi}{e^2q_nn}\ge\frac{16\varpi}{625e^2n}.
\]
\end{corollary}

\begin{proof}
Translation invariance and reflection positivity are identities and
inequalities among finitely many local functions of $\mu_L$~\cite{OSei78}, and
pass to $\omega$. For a continuous gauge-invariant $F$ at positive times the
supports of $\vartheta(\bar F)$ and $\tau_tF$ are at distance $d$ with $|d-t|\le
c_F$, and Theorem~\ref{thm:decay} on $\Lambda_L$ gives $|\Cov_{\mu_L}(\vartheta(\bar
F),\tau_tF)|\le C_F(1+t)e^{-m_nt}$ uniformly in $L$. In the limit
$|\omega(\vartheta(\bar F)\tau_tF)-|\omega(F)|^2|\le C_F(1+t)e^{-m_nt}$, and
Proposition~\ref{prop:gapdecay}, applied with every $m<m_n$, gives the gap.
\end{proof}

\subsection{The continuum limit}

Let $a_k\to0$ be lattice spacings with $a_k=2^{-k}a_0$, and $\beta_k$ couplings.
Let $\omega_k$ be lattice states with transfer matrices $\mathbb T_k$ on spaces
$\mathcal H_k$ with vacua $\Omega_k$. Let $\mathcal T$ be a complex vector space of
test functions on $\R^4$ supported in $\{t>0\}$, and for each $k$ let $f\mapsto
O_k(f)$ be a linear map to gauge-invariant local observables at positive
times; the example is $O_k(f)=Z_k\sum_xa_k^4f(a_kx)\Phi(x)$ with a local field
$\Phi$ and a renormalisation constant $Z_k$.

\begin{definition}\label{def:limit}
A \emph{continuum limit} of $(\omega_k,O_k)$ is a Hilbert space $\mathcal H$ with
a unit vector $\Omega$, a self-adjoint $H\ge0$ with $H\Omega=0$, and a linear
map $\psi:\mathcal T\to\mathcal H$ with dense range such that for all
$f,g\in\mathcal T$ and dyadic $\tau\ge0$
\[
\ip{\psi(f)}{e^{-\tau H}\psi(g)}=\lim_k\omega_k\bigl(\vartheta(\overline{O_k(f)})
\,\tau_{\tau/a_k}O_k(g)\bigr),\qquad
\ip\Omega{\psi(f)}=\lim_k\omega_k(O_k(f)),
\]
over the $k$ with $\tau/a_k\in\Z$. It is \emph{nontrivial} if $\mathcal H\ne\C
\Omega$.
\end{definition}

The Osterwalder--Schrader reconstruction produces such a limit from
convergent renormalised Schwinger functions satisfying the
axioms~\cite{OS73,OS75}; the definition records only what is used below.

\begin{proof}[Proof of Theorem~\ref{thm:B}]
Let $\omega_k$ be torus limits at $\beta_k$ along the tori on which
$\BP(n_k,\varpi)$ holds. By Corollary~\ref{cor:latticegap}, $\mathbb T_k$ has gap
at least $16\varpi/(625e^2n_k)$ on $\Omega_k^\perp$. For $f\in\mathcal T$ put
$O=O_k(f)$ and $\varphi_k=\Psi_O-\omega_k(O)\Omega_k$. For dyadic $\tau$ and $k$
large, $t=\tau/a_k$ is an integer and $16\varpi t/(625e^2n_k)\ge\Delta\tau$ since
$n_ka_k\le\ell$, so
\[
\omega_k\bigl(\vartheta(\bar O)\tau_tO\bigr)-|\omega_k(O)|^2=\ip{\varphi_k}{\mathbb T_k^t
\varphi_k}\le e^{-\Delta\tau}\|\varphi_k\|^2 .
\]
Both sides converge by Definition~\ref{def:limit}: with $\varphi=\psi(f)-\ip\Omega
{\psi(f)}\Omega$, $\ip\varphi{e^{-\tau H}\varphi}\le e^{-\Delta\tau}\|\varphi\|^2$. These
vectors form a subspace dense in $\Omega^\perp$, and the inequality extends to
$\Omega^\perp$ and to all $\tau\ge0$ by continuity. For $\varphi\in\Omega^\perp$ with
spectral measure $\mu_\varphi$ of $H$ this says $\int e^{-\tau E}d\mu_\varphi(E)\le e^{-\Delta
\tau}\mu_\varphi([0,\infty))$ for all $\tau$, which forces $\mu_\varphi([0,\Delta-\eta])=0$
for every $\eta>0$. So $0$ is a simple eigenvalue and the spectrum on
$\Omega^\perp$ lies in $[\Delta,\infty)$, and is not empty if $\Omega^\perp\ne\{0\}$.
\end{proof}

The proof uses the cube condition only through the decay of two-point
functions of gauge-invariant observables, and the constant $\Delta$ depends on
$\varpi$ and $\ell$ alone.
